\section*{前書き}
データベースという言葉を知ったのは，大学3年生の授業でした．
当時，工学部情報科学科の学生であった私は，必修科目である「データベース」の講義を受けていましたが，他の授業と同様，ただ漫然と講義を受けているだけでした．
単位を修得したものの，データベースの意義をまったく理解していませんでした．

その後，縁があってデータベース研究者として名を馳せた先生の研究室に入りましたが，研究室の中心テーマは情報検索であり，研究や仕事でデータベースを使うことはあっても，データベースの理論は遠い遠い存在でした．
そんな自分がまさかデータベースの教科書を書くことになるとは，夢にも思っていませんでした．

さて，本書は名古屋市立大学データサイエンス学部2年生の授業科目「データベース」のために作成した講義ノートをもとに書かれています．
授業を設計する際，GoogleやAmazonなど，巨大IT企業がクラウド上で高性能なデータベースを提供してくれる今日，データベースを専門的に教える必要性があるのだろうかと悩みました．
ましてや教える相手は，コンピュータサイエンスを専門としないデータサイエンス学部の学生です．
今のご時世，データベースの授業で何を教えればよいのでしょうか．

データ活用が叫ばれる中，私たちの周りには，このままでは管理も分析もままならない，目も当てられないようなデータが大量に溢れています．
大学で何かを専門的に学んだ研究者であっても，データ分析にはExcelを使っていて，いざ大量のデータが出てきたときに何もできずに困っている人が多くいます．
データと上手く付き合っていくためには，データを正しく管理し，大量のデータから欲しい情報を見つけるための「基礎」だけでも知っておくことは悪くありません．
むしろ，必須と言えるでしょう．
データにまつわる悲劇を見ているうちに，そう強く思うようになりました．

そこで本書では，データベースという概念の中核かつ基礎である「関係データベース」を対象に，「設計」と「問い合わせ（SQL）」の2点に焦点を絞り，その基礎について解説を行っています．
したがって，関係データベースの教科書では鉄板ネタである以下のようなトピックはあえて扱っていません．
\begin{itemize}
    \item 関係論理，関係代数
    \item トランザクション
    \item クエリ最適化
    \item 物理的なデータの格納方式
    \item 障害回復
\end{itemize}

とはいえ，世の中にはデータベースに関する書籍が溢れています．
和書の教科書であれば，データベース分野で著名な増永先生，北川先生，吉川先生がそれぞれ名著を書いておられます．
技術書であれば，実践的で最先端の話題を取り上げた本がたくさん出ていますし，最近ではYouTubeに分かりやすい解説動画も多数アップロードされています．
そこで本書を執筆するにあたっては，学術的な雰囲気を保ちつつも，データ管理やデータ分析の経験に乏しい初学者でも分かりやすい書籍づくりを心がけました．
本書が，先に挙げたような，より専門的な教科書や実践的な技術書への橋渡しになれば幸いです．

本書を著すにあたっては，多くの方々のお世話になりました．
兵庫県立大学社会情報科学部の大島裕明先生には，草稿だけでなく出発点となった講義資料づくりにも有益なフィードバックをいただきました．
また，本書の執筆を勧めてくださった森北出版の福島さんには，「重鎮の先生方が教科書を書いているのに，私なんぞが教科書を書くなんて，私の立場も理解してくださいよ」と何度か泣き言を言いましたが，その度に叱咤激励を頂戴しました．
ここに厚く御礼申し上げます．
最後に，土日にも嫌な顔をせずに執筆時間を与えてくれた妻に深く感謝いたします．

\vspace{2zw}
\begin{flushright}
山本 祐輔
\end{flushright}
